% !TEX root = ./demo.tex
\documentclass[
  lang=english,
  accent=0055A4,      % Royal Academic Navy Blue
  circuits=true,      % Native Circuitikz support
  qed=smiley          % Classic smiley QED tombstone
]{opalith}

\title{Digital Transmission and Embedded Systems}
\subtitle{From Theoretical Signal Foundations to Hardware Conditioning and Real-Time Control}
\author{Simone Lapomarda}
\professor{Prof. A. Rossi}
\institution{University of Padua}
\department{Department of Information Engineering}
\degree{BSc in Computer Engineering}
\coursecode{Course ID: IN01123456 \quad$\cdot$\quad 9 CFU}
\academicyear{2026/2027}
\headertitle{Digital Transmission and Embedded Systems}

\begin{document}

\maketitle
\tableofcontents

% ==============================================================================
\chapter{Foundations of Signal Theory and Sampling}
% ==============================================================================

\section{Bandlimited Signals and the Frequency Domain}
\label{sec:signals_frequencies}

In \uline{digital communication systems}, the fundamental goal is to transmit binary information across a physical communication channel subject to finite bandwidth constraints, linear phase distortion, and additive thermal noise.

Consider a deterministic continuous-time signal $x(t) \in \mathcal{L}^2(\RR)$. Its \textbf{Fourier transform} is conventionally defined as:
\begin{equation}
  X(f) = \mathcal{F}\{x(t)\} = \int_{-\infty}^{+\infty} x(t)\, e^{-j 2\pi f t}\, \dd t
  \label{eq:fourier_def}
\end{equation}
with the corresponding inverse synthesis transform given by:
\begin{equation}
  x(t) = \mathcal{F}^{-1}\{X(f)\} = \int_{-\infty}^{+\infty} X(f)\, e^{j 2\pi f t}\, \dd f
\end{equation}

\dfn{Bandlimited Signal}{
  A continuous-time signal $x(t)$ is said to be \uline{strictly bandlimited} in the Fourier sense if there exists a finite cut-off frequency $B > 0$, termed the \textbf{nominal bandwidth}, such that its spectral energy density vanishes identically outside $[-B, B]$:
  \[
    X(f) = 0, \quad \forall |f| > B
  \]
}

By Parseval's energy conservation theorem, the total energy $E_x$ of a finite-energy signal is preserved identically across domains:
\begin{equation}
  E_x = \int_{-\infty}^{+\infty} |x(t)|^2 \, \dd t = \int_{-\infty}^{+\infty} |X(f)|^2 \, \dd f = \int_{-B}^{+B} |X(f)|^2 \, \dd f
\end{equation}
In practical engineering systems, strict bandlimiting ensures that $x(t)$ is smooth and analytic over $\RR$, precluding sudden jump discontinuities in the time domain and guaranteeing strong temporal correlation between adjacent samples.

\section{The Nyquist-Shannon Sampling Theorem}
\label{sec:nyquist_shannon}

Sampling represents the fundamental bridge between the continuous physical realm and discrete digital processing. In periodic uniform sampling, a continuous signal $x(t)$ is evaluated at fixed temporal intervals $T_s$:
\begin{equation}
  x[n] \defeq x(n T_s), \quad f_s \defeq \frac{1}{T_s}
\end{equation}
where $f_s$ is the sampling rate in Hertz. If the conversion process is performed too slowly, high-frequency spectral components alias irreversibly into the baseband, permanently destroying signal fidelity. The central question answered by Harry Nyquist (1928) and formalized by Claude Shannon (1949) is: \textit{what is the minimal sampling frequency required to capture continuous information without irreversible degradation?}

\newpage
\thm{Sampling Theorem (Nyquist-Shannon)}{
  \label{thm:nyquist}
  Let $x(t)$ be a continuous-time signal strictly bandlimited to $[-B, B]$. If the sampling frequency $f_s$ satisfies the non-overlapping condition:
  \begin{equation}
    f_s \ge 2B
    \label{eq:nyquist_rate}
  \end{equation}
  then the continuous signal $x(t)$ can be reconstructed \textbf{exactly and without distortion} for all $t \in \RR$ from its discrete samples $x(n T_s)$ via cardinal sinc interpolation:
  \begin{equation}
    x(t) = \sum_{n=-\infty}^{+\infty} x(n T_s) \operatorname{sinc}\left( \frac{t - n T_s}{T_s} \right)
  \end{equation}
  where $\operatorname{sinc}(u) \defeq \frac{\sin(\pi u)}{\pi u}$ for $u \ne 0$ and $\operatorname{sinc}(0) = 1$.
}

\dem{Spectral Proof}{
  Ideal impulse sampling can be modeled as the multiplication of $x(t)$ by a periodic Dirac comb (Shah distribution) $p(t) = \sum_{n=-\infty}^{+\infty} \delta(t - n T_s)$:
  \[
    x_p(t) = x(t) \cdot \sum_{n=-\infty}^{+\infty} \delta(t - n T_s) = \sum_{n=-\infty}^{+\infty} x(n T_s) \delta(t - n T_s)
  \]
  Applying the Fourier property that maps a time-domain product to a spectral convolution, and recalling that $\mathcal{F}\{p(t)\} = f_s \sum_{k=-\infty}^{+\infty} \delta(f - k f_s)$:
  \[
    X_p(f) = X(f) * P(f) = f_s \sum_{k=-\infty}^{+\infty} X(f - k f_s)
  \]
  The sampled spectrum $X_p(f)$ consists of an infinite summation of replicas centered at integer multiples $k f_s$. When $f_s - B \ge B \iff f_s \ge 2B$, adjacent spectral replicas do not overlap. The baseband replica ($k=0$) can be isolated by applying an ideal low-pass brickwall filter:
  \[
    H_{rec}(f) = T_s \, \operatorname{rect}\left(\frac{f}{2B}\right) = \begin{cases} T_s & |f| \le B \\ 0 & |f| > B \end{cases}
  \]
  Multiplying $X(f) = X_p(f) \cdot H_{rec}(f)$ and transforming back to the time domain yields the classical Whittaker-Shannon cardinal series.
}

\begin{figure}[htbp]
\centering
\begin{tikzpicture}
\begin{axis}[
    width=0.88\textwidth, height=5.2cm,
    axis lines=middle,
    axis on top=true,
    xlabel={$f$ [kHz]},
    xlabel style={at={(ticklabel* cs:1.02)}, anchor=west, xshift=2.5mm},
    ylabel={$X_p(f)$},
    ylabel style={at={(ticklabel* cs:1.02)}, anchor=south, yshift=1mm},
    xmin=-9, xmax=9.8,
    ymin=0, ymax=1.45,
    xtick={-6, -4, -2, 0, 2, 4, 6},
    xticklabels={$-f_s$, , $-B$, $0$, $B$, , $f_s$},
    ytick={1},
    yticklabels={$f_s X_0$},
    yticklabel style={anchor=south east, font=\small, xshift=-1.5mm, yshift=1mm, fill=white, fill opacity=0.9, text opacity=1, inner sep=1.5pt},
    every axis plot/.append style={thick},
    legend style={at={(0.98,0.92)}, anchor=north east, font=\small}
]
  % Central replica (k=0)
  \addplot[color=doc, fill=doc!15] coordinates {(-2,0) (-1,1) (0,1) (1,1) (2,0)} \closedcycle;
  \addlegendentry{Baseband spectrum ($k=0$)}

  % Right replica (k=+1, fs = 6 kHz)
  \addplot[color=blue!80!black, fill=blue!10] coordinates {(4,0) (5,1) (6,1) (7,1) (8,0)} \closedcycle;
  \addlegendentry{Periodic replicas ($k = \pm 1$)}

  % Left replica (k=-1, fs = -6 kHz)
  \addplot[color=blue!80!black, fill=blue!10] coordinates {(-8,0) (-7,1) (-6,1) (-5,1) (-4,0)} \closedcycle;

  % Reconstruction filter
  \addplot[color=red, dashed, line width=1.1pt] coordinates {(-3,0) (-3,1.2) (3,1.2) (3,0)};
  \addlegendentry{Reconstruction filter $H(f)$}
\end{axis}
\end{tikzpicture}
\caption{Frequency-domain representation of periodic sampling without aliasing ($f_s > 2B$).}
\label{fig:sampling_spectrum}
\end{figure}

\newpage
\section{Inter-Symbol Interference and Nyquist's Criterion}
\label{sec:isi_nyquist}

When modulated pulses propagate across a bandlimited channel, dispersive spreading causes adjacent pulse tails to overlap, generating \uline{Inter-Symbol Interference} (ISI).

\frm{Nyquist Criterion for Zero ISI}{
  To ensure that a pulse-shaping filter $g(t)$ causes no interference at optimum symbol decision instants $t = m T$, its frequency response $G(f)$ must satisfy:
  \begin{equation}
    \sum_{k=-\infty}^{+\infty} G\left(f - \frac{k}{T}\right) = T, \quad \forall |f| \le \frac{1}{2T}
    \label{eq:nyquist_isi_criterion}
  \end{equation}
}

\oss{Spectral Efficiency vs. Jitter Sensitivity}{
  The theoretical minimum-bandwidth pulse satisfying \eqref{eq:nyquist_isi_criterion} is the sinc pulse, achieving $B_{min} = \frac{1}{2T} = \frac{R_s}{2}\text{ Hz}$. However, because sinc pulses decay as slowly as $1/t$, small timing offsets or clock jitter induce divergent cumulative ISI series.
}

\ex{Raised Cosine Spectrum}{
  To ensure rapid temporal decay without infinite tails, communication systems employ the \textbf{raised cosine} filter with roll-off factor $\alpha \in [0, 1]$. The total transmission bandwidth is given by:
  \[
    B = \frac{1 + \alpha}{2T} = \frac{R_s}{2}(1 + \alpha)
  \]
  For a link operating at $R_s = 100\text{ kbaud}$ with roll-off $\alpha = 0.25$, the required channel bandwidth is $B = \frac{100}{2}(1.25) = 62.5\text{ kHz}$.
}

% ==============================================================================
\chapter[Analog Conditioning and Front-End]{Analog Conditioning\\and Electronic Front-End}
% ==============================================================================

\section{Active Anti-Aliasing Filter Design}
\label{sec:active_filter}

Before an analog signal can be safely converted by an ADC, an analog \uline{anti-aliasing filter} must attenuate spectral noise and high-frequency interferers above the Nyquist frequency $f_s/2$.

\begin{figure}[htbp]
\centering
\begin{circuitikz}[european, scale=0.82, transform shape]
    % Op Amp
    \draw (4.5, 0) node[op amp] (opamp) {};
    \draw (opamp.+) node[ground] {};

    % Input branch
    \draw (0, 0.5) node[left] {$v_{in}(t)$}
          to[R, l=$R_1$, o-*] (2, 0.5) -- (opamp.-);

    % Feedback branch R2
    \draw (2, 0.5) -- (2, 2.0)
          to[R, l=$R_2$] (6, 2.0) -- (6, 0)
          to[short, *-o] (7.2, 0) node[right] {$v_{out}(t)$};

    % Parallel capacitor branch
    \draw (2, 2.0) -- (2, 3.2)
          to[C, l=$C_1$] (6, 3.2) -- (6, 2.0);

    % Output
    \draw (opamp.out) -- (6, 0);
\end{circuitikz}
\caption{First-order inverting active low-pass filter stage for ADC interface.}
\label{fig:active_lowpass}
\end{figure}

Nodal analysis at the inverting virtual ground yields the Laplace transfer function:
\begin{equation}
  W(s) = \frac{V_{out}(s)}{V_{in}(s)} = -\frac{Z_2(s)}{Z_1(s)} = -\frac{R_2 \parallel \frac{1}{sC_1}}{R_1} = -\frac{R_2}{R_1} \cdot \frac{1}{1 + s R_2 C_1}
\end{equation}

\prp{Cut-Off Frequency and Passband Gain}{
  Defining $A_0 \defeq -\frac{R_2}{R_1}$ (DC passband gain) and $\omega_H \defeq \frac{1}{R_2 C_1}$, the frequency response magnitude is:
  \begin{equation}
    |W(j\omega)| = \frac{|A_0|}{\sqrt{1 + (\omega / \omega_H)^2}}
  \end{equation}
  At the cut-off angular frequency $\omega = \omega_H$, the gain drops by $\sqrt{2}$, corresponding to the conventional $-3\text{ dB}$ attenuation threshold.
}

\warn{Saturation and Supply Rails}{
  Real operational amplifiers cannot swing beyond their power supply voltages:
  \[
    V_{EE} + V_{sat} \le v_{out}(t) \le V_{CC} - V_{sat}
  \]
  Driving the amplifier beyond these limits induces severe \uline{harmonic clipping}, producing spurious high-frequency harmonics that fold back into the baseband spectrum as in-band aliasing.
}

\nt{Single-Supply DC Biasing}{
  In modern unipolar embedded systems ($0 - 3.3\text{ V}$ microcontrollers), the non-inverting terminal must be biased at $V_{ref} = V_{DD}/2 = 1.65\text{ V}$ via a decoupled resistive voltage divider to center the dynamic range.
}

\section{Placeholder for High-Frequency Front-End Schematic}

Lecture slides and hardware datasheets often contain complex multi-stage diagrams that must be cropped and integrated during revision:

\begin{center}
  \figplaceholder{Slide 22 (Prof. Rossi): Superheterodyne receiver front-end block diagram with I/Q demodulation mixer and Programmable Gain Amplifier (PGA)}
\end{center}

\section{Practical Component Sizing Exercise}

\exer{Anti-Aliasing Filter Component Sizing}{
  Design the active first-order low-pass filter shown in Figure~\ref{fig:active_lowpass} to interface a seismic sensor to an ADC running at $f_s = 20\text{ kHz}$.
  The design specifications are:
  \begin{enumerate}[label=\alph*.]
    \item Half-power bandwidth ($-3\text{ dB}$ cut-off) at $f_H = 4\text{ kHz}$.
    \item In-band passband voltage gain magnitude $|A_0| = 5$ ($14\text{ dB}$).
    \item Input resistance $R_1 \ge 10\text{ k}\Omega$ to prevent loading effects on the transducer.
  \end{enumerate}
  Select standard commercial values (E12 series) for $R_1$, $R_2$, and $C_1$.
}

\begin{Solution}
\begin{enumerate}[label=\alph*.]
  \item \textbf{Selection of $R_1$ and $R_2$}:
  Setting the input resistance to $R_1 = 10\text{ k}\Omega$ to satisfy the source loading constraint, the feedback resistance required for a gain $|A_0| = R_2/R_1 = 5$ is:
  \[
    R_2 = 5 \cdot 10\text{ k}\Omega = 50\text{ k}\Omega \quad \xrightarrow{\text{E12 series}} \quad R_2 = 47\text{ k}\Omega
  \]
  giving an effective passband gain $|A_{0,eff}| = 4.7$ (less than $6\%$ design deviation).

  \item \textbf{Capacitance Calculation for $C_1$}:
  From the cut-off frequency relationship $\omega_H = 2\pi f_H = 1/(R_2 C_1)$:
  \[
    C_1 = \frac{1}{2\pi f_H R_2} = \frac{1}{2\pi \cdot 4000 \cdot 47 \times 10^3} \approx 846\text{ pF} \quad \xrightarrow{\text{E12 series}} \quad C_1 = 820\text{ pF}
  \]
\end{enumerate}
\end{Solution}

% ==============================================================================
\chapter{Digital Processing and Embedded Real-Time Implementation}
% ==============================================================================

\section{Circular Buffer Acquisition Algorithm}
\label{sec:circular_buffer}

In real-time microcontroller architectures (e.g., ARM Cortex-M or RISC-V), ADC conversions are triggered by periodic hardware timers and transferred into memory via DMA or fast interrupts. To decouple acquisition rates from computation deadlines without introducing \hyper{sec:nyquist_shannon}{data races}, a \uline{lock-free circular ring buffer} is utilized.

\begin{algorithm}[H]
\caption{Real-time causal FIR filtering on a circular ring buffer.}
\label{algo:fir_buffer}
\KwIn{New ADC conversion sample x_new, write index pointer w_idx}
\KwOut{Updated buffer B[0...N-1], computed filter output y}
\BlankLine
B[w_idx] \gets x_new\;
w_idx \gets (w_idx + 1) \bmod N\;
acc \gets 0\;
\For{k \gets 0 \KwTo M - 1}{
  idx \gets (w_idx - 1 - k + N) \bmod N\;
  acc \gets acc + h[k] * B[idx]\;
}
\Return{acc}\;
\end{algorithm}

\section{Embedded C Implementation}
\label{sec:c_code}

The listing below illustrates the production C Interrupt Service Routine (ISR) executing on a 32-bit embedded target:

\begin{lstlisting}[caption={Embedded C timer ISR for periodic ADC sampling and FIR filtering.}, label={lst:fir_isr}]
#include <stdint.h>
#define BUFFER_SIZE   64
#define TAPS_COUNT    8

static volatile uint16_t adc_buffer[BUFFER_SIZE];
static volatile uint8_t  head_idx = 0;
// Symmetric FIR coefficients (fixed-point Q15 representation)
static const int16_t fir_coeffs[TAPS_COUNT] = {
    1024, 3450, 7812, 11420, 11420, 7812, 3450, 1024
};

void ADC_Timer_IRQHandler(void) {
    // 1. Hardware sample acquisition & circular buffer insert
    adc_buffer[head_idx] = ADC1->DR;
    head_idx = (head_idx + 1) & (BUFFER_SIZE - 1); // Fast power-of-two wrap

    // 2. Compute discrete convolution sum
    int32_t accumulator = 0;
    for (int k = 0; k < TAPS_COUNT; ++k) {
        uint8_t sample_idx = (head_idx - 1 - k) & (BUFFER_SIZE - 1);
        accumulator += (int32_t)fir_coeffs[k] * adc_buffer[sample_idx];
    }

    // 3. Output scaling to unipolar DAC range
    int16_t filtered = (int16_t)(accumulator >> 15);
    DAC1->DHR12R1 = (uint16_t)(filtered + 2048);
}
\end{lstlisting}

\section{Comparative Filter Performance Benchmark}
\label{sec:benchmark_table}

Table~\ref{tab:filter_comparison} summarizes the computational trade-offs measured on an ARM Cortex-M4 microcontroller running at $168\text{ MHz}$ across digital filter topologies.

\begin{table}[htbp]
\centering
\begin{tabular}{lcccc}
\toprule
\textbf{Filter Architecture} & \textbf{Order} & \textbf{CPU Cycles/Sample} & \textbf{Latency [$\mu$s]} & \textbf{Stability} \\
\midrule
Simple Moving Average (SMA)  & 8               & 18                         & 0.11                      & Unconditional      \\
Raised Cosine FIR            & 16              & 42                         & 0.25                      & Unconditional      \\
Raised Cosine FIR            & 32              & 85                         & 0.51                      & Unconditional      \\
Butterworth IIR (Biquad)     & 4 (2 biquads)   & 24                         & 0.14                      & Conditional        \\
Chebyshev Type I IIR         & 4 (2 biquads)   & 26                         & 0.15                      & Conditional        \\
\bottomrule
\end{tabular}
\caption{Complexity and execution latency across digital filter topologies.}
\label{tab:filter_comparison}
\end{table}

\section{Real-Time Schedulability Verification Exercise}

\exer{Real-Time Deadline Schedulability Analysis}{
  Consider the embedded ISR routine defined in Listing~\ref{lst:fir_isr} executing on a processor clocked at $f_{cpu} = 48\text{ MHz}$.
  Given the following system constraints:
  \begin{itemize}
    \item The sampling frequency is set to $f_s = 50\text{ kHz}$.
    \item Computational filtering requires $140$ core CPU execution cycles.
    \item Hardware interrupt context-switching overhead (prologue and epilogue) is $24$ cycles.
  \end{itemize}
  Compute the \uline{CPU utilization factor} ($U$) and determine whether the hard real-time deadline is strictly met.
}

\begin{Solution}
The total period available between successive timer interrupts is:
\[
  T_s = \frac{1}{f_s} = \frac{1}{50\,000\text{ Hz}} = 20\,\mu\text{s}
\]
The total cycle count consumed by the processor per sample is:
\[
  C_{tot} = C_{calc} + C_{overhead} = 140 + 24 = 164 \text{ cycles}
\]
The physical execution time is:
\[
  T_{exe} = \frac{C_{tot}}{f_{cpu}} = \frac{164}{48 \times 10^6\text{ s}^{-1}} \approx 3.417\,\mu\text{s}
\]
The CPU utilization factor is given by:
\[
  U = \frac{T_{exe}}{T_s} = \frac{3.417\,\mu\text{s}}{20\,\mu\text{s}} \approx 0.1708 \implies 17.08\%
\]
Because $U < 100\%$ (and well below the classical Liu-Layland real-time bound of $69\%$), the system is \textbf{fully schedulable}, leaving more than $82\%$ of CPU bandwidth available for background application processing.
\end{Solution}

\end{document}
